\chapter{このテキストの使い方}

\section{誰のためのテキストか}

このテキストは、次のような読者を想定している。

\begin{itemize}
\item Python をこれから学びたい人
\item 多少は書けるが、解析コードをどう整理すればよいか分からない人
\item 小さなデータでは動くが、大きなデータになると急に遅くなる理由を知りたい人
\item 結果の図をそれらしく描くことはできても、正しさの確認やフィットの説明に自信がない人
\end{itemize}

逆に、すでに Python で大規模解析を日常的に行っており、NumPy や pandas の内部挙動にも慣れている読者にとっては、前半の内容は基礎的に見えるだろう。その場合は、後半の高速化や総合演習から読んでも構わない。

本書は、単に「使い方を覚える」ためだけの資料ではない。研究や課題で自力の解析を組み立てるとき、何を順に確かめればよいか、どの段階で道具を変えるべきかまで含めて学ぶことを目指している。したがって、短いサンプルコードであっても、なぜその書き方にしているのかを説明する。

\section{何を順に学ぶのか}

本書では、大きく分けて三段階で学ぶ。

第一段階では、Python の基本的な書き方と、ファイルを読むための最低限の技術を扱う。ここでは、短いスクリプトを自分で動かし、何が入力で何が出力なのかをはっきり意識する。

第二段階では、NumPy、pandas、Polars、Matplotlib、SciPy、Astropy を使い、実際の解析に必要な表現力を身につける。配列、表形式データ、ヒストグラム、フィット、時刻と単位の扱いが、この段階の中心になる。

第三段階では、高速化と設計を学ぶ。遅い実装と速い実装を比べ、どこで時間を使っているのかを測り、アルゴリズムの違いが速度にどう響くのかを考える。この延長として、Python だけでなく C や C++ へ処理を移す発想も扱う。

この順序には理由がある。NumPy や pandas を早く使いたくなるのは自然だが、入力と出力、型、反復、条件分岐が曖昧なままでは、便利なライブラリを使っても何が起きているか分からなくなる。基礎を先に固めるのは遠回りではなく、後半を理解する最短経路である。

\section{本書で大切にしたい姿勢}

本書で繰り返し強調するのは、次の三点である。

\begin{itemize}
\item 結果をそのまま信じず、自分で確かめること
\item 遅いと感じたら、どこが遅いかを測ること
\item コードが動く理由と、速くなる理由を言葉で説明すること
\end{itemize}

特にデータ解析では、コードが実行できたことと、結果が正しいことは同じではない。さらに、正しいコードであっても、データ量が増えた途端に実用にならなくなることがある。本書では、その両方を区別して考える。

もう一つ大切なのは、結果を言葉に直すことである。たとえば「NumPy を使ったら速くなった」と書くだけでは不十分で、どの操作を Python ループから配列演算へ移したのか、どの程度速くなったのか、メモリ使用量はどう変わったのか、といった説明が必要になる。レポートや口頭説明では、この差がそのまま理解の差として現れる。

\section{章の読み進め方}

各章には、できるだけ次の順序で内容を置く。

\begin{enumerate}
\item その章で扱う考え方の意味
\item 最小限のコード例
\item よくある落とし穴
\item 解析へのつながり
\end{enumerate}

最初は、コードを丸ごと理解しようとしなくてよい。まずは「何を入力し、何を出しているか」を追う。その後で「なぜこの書き方を選んだか」を読むと、理解しやすい。

特に初学者は、1 行ずつ読んで意味を取ろうとして全体像を見失いがちである。まずは関数名、引数、返り値、ファイル名、出力先だけを追い、流れをつかんでから細部へ戻る方が理解しやすい。本書のコード例も、その読み方を意識して配置している。

\section{併用するとよい公開資料}

本書は独立して読めるように書いているが、コマンドの細かいオプションやライブラリの仕様は版によって変わることがある。そのため、必要に応じて公式の公開資料も併用した方がよい。本書の記述も、以下のような公開資料の内容を踏まえて整理している。

\begin{itemize}
\item Python 全般: \href{https://docs.python.org/3/tutorial/}{Python 公式チュートリアル} と \href{https://docs.python.org/3/library/}{Python 標準ライブラリ}
\item JupyterLab: \href{https://jupyterlab.readthedocs.io/}{JupyterLab Documentation}
\item uv: \href{https://docs.astral.sh/uv/}{uv Documentation}
\item pyenv: \href{https://github.com/pyenv/pyenv}{pyenv README / Documentation}
\item NumPy: \href{https://numpy.org/doc/stable/}{NumPy Documentation}
\item pandas: \href{https://pandas.pydata.org/docs/}{pandas Documentation}
\item Polars: \href{https://docs.pola.rs/}{Polars Documentation}
\item Matplotlib: \href{https://matplotlib.org/stable/}{Matplotlib Documentation}
\item SciPy: \href{https://docs.scipy.org/doc/scipy/}{SciPy Documentation}
\item Astropy: \href{https://docs.astropy.org/}{Astropy Documentation}
\item Git: \href{https://git-scm.com/doc}{Git Documentation}
\item GitHub: \href{https://docs.github.com/}{GitHub Docs}
\end{itemize}

特にシェルの \code{ssh}、\code{rsync}、\code{ls}、\code{chmod} などは、環境ごとに細かな違いがある。細部を確認したいときは \code{man ssh} や \code{man rsync} のようにマニュアルを開く習慣も有効である。

\section{この本の後半で扱う実践的な内容}

本書の後半では、単なる文法紹介ではなく、実際のデータを扱うときに何が問題になるかを扱う。そこでは、正しさの確認、バグの切り分け、図の作成、フィット、速度改善といった要素が一つの流れとして現れる。

最後の章には具体的な総合演習を置くが、それは本書で扱った一般的な考え方をまとめて使うための実践例という位置付けである。したがって、前半の章では特定の演習に閉じず、より広いデータ解析へそのまま応用できる書き方と考え方を重視する。

言い換えると、後半へ行くほど「何を書くか」よりも「どう考えるか」が主題になる。ライブラリの関数名は将来変わることがあるが、入力を確かめる、分布を見る、速度を測る、遅い理由を切り分ける、といった考え方は長く残る。
